\chapter{Walkthrough: Play the Whole Book Through Two Short Songs}
\label{ch:walkthrough}

\begin{quote}
\itshape Arrangement is not using each chapter once; it is letting them all take their place in one song. A walkthrough goes section by section, puts each layer in, and checks that it does not steal another layer's seat.
\end{quote}

\section{The goal: two short songs, one workflow}
\label{sec:walkthrough-goal}

This book ends without a ``universal recipe''. Instead it uses two \term{short songs}{short-song} as \term{demo projects}{demo-project}: one folk, one pop. Neither is a commercial length; each is sixteen bars (eight-bar verse + eight-bar chorus), but both follow the same workflow:

\begin{enumerate}
\item Write the section chart and energy curve;
\item Create and name tracks; assign the four layer roles;
\item Build the rhythm section, then the harmony layer;
\item Write the lead melody/vocal;
\item Add colour layers (strings, woodwinds/brass, electronic);
\item Walk through section by section and fix seat-stealing;
\item Save the project;
\item Export WAV and stems, and write a delivery receipt.
\end{enumerate}

Neither demo is a preinstalled project in the repository; both are teaching setups built by the methods in this book. You can rebuild them in Yeban from the tables, or use them only as a checklist.

\section{Demo one: the folk short song ``Window Side''}
\label{sec:folk-demo}

The folk principle is ``few but clear'': lead vocal, acoustic guitar, bass, light drums/percussion, string pad. The sixteen-bar section chart:

\begin{longtable}{p{0.18\textwidth}p{0.20\textwidth}p{0.52\textwidth}}
\toprule
Bars & Section & Arrangement action \\
\midrule
\endfirsthead
\toprule
Bars & Section & Arrangement action \\
\midrule
\endhead
1--8 & Verse & Acoustic guitar + bass + light drums; lead; strings stay out \\
9--12 & First chorus & Add string pad and full drums; lead stays centred \\
13--16 & Second chorus & Add one double/harmony; strings hold long notes, no counterpoint \\
\bottomrule
\end{longtable}

\begin{longtable}{p{0.20\textwidth}p{0.18\textwidth}p{0.52\textwidth}}
\toprule
Track & Type & Role and writing \\
\midrule
\endfirsthead
\toprule
Track & Type & Role and writing \\
\midrule
\endhead
Lead vocal & Audio & Narrative; leave one beat at phrase ends, do not fill every gap \\
Acoustic guitar & Audio & Harmonic skeleton; trim clips to leave the vocal's phrase ends \\
Bass & MIDI/audio & Root; mostly same-note repetition, few leaps \\
Drums/percussion & MIDI/audio & Light skeleton; verse keeps kick/snare/light tambourine \\
String pad & MIDI & Enters in the chorus; long notes, not the vocal's rhythm \\
Double/harmony & Audio & Enters after bar 13; low third, not in the lead's register \\
\bottomrule
\end{longtable}

The folk energy curve rises gently; do not turn the chorus into electronic dance music. The chorus is ``bigger'' because of the string pad, fuller drums and one double, not because of more tracks.

\section{Demo two: the pop short song ``Night Run''}
\label{sec:pop-demo}

The pop principle is ``each layer clearer'', not ``each layer more''. Sixteen bars:

\begin{longtable}{p{0.18\textwidth}p{0.20\textwidth}p{0.52\textwidth}}
\toprule
Bars & Section & Arrangement action \\
\midrule
\endfirsthead
\toprule
Bars & Section & Arrangement action \\
\midrule
\endhead
1--8 & Verse & Drums + bass + pad + lead; pluck only at phrase ends \\
9--12 & First chorus & Add double, continuous pluck, brass stab at bar 9 \\
13--16 & Second chorus & Add harmony and stack in the last two bars; brass stops after the stab \\
\bottomrule
\end{longtable}

\begin{longtable}{p{0.20\textwidth}p{0.18\textwidth}p{0.52\textwidth}}
\toprule
Track & Type & Role and writing \\
\midrule
\endfirsthead
\toprule
Track & Type & Role and writing \\
\midrule
\endhead
Lead vocal & Audio & Lyric first; lift in the chorus but stay in a comfortable register \\
Double & Audio & Same melody, second take; not a copy of the lead \\
Harmony & Audio & A third below the lead; chorus only \\
Drums & MIDI & Four-beat skeleton; add density in the chorus \\
Bass & MIDI & Root; divide the low end with sub \\
Pad & MIDI & Midrange bed; open cutoff in the chorus \\
Pluck & MIDI & High decoration; phrase ends/gaps, not continuous \\
Brass & MIDI & Short stab at bar 9; exit after a beat and a half \\
\bottomrule
\end{longtable}

The pop chorus is not ``open every track''; it is ``each layer more complete'': fuller drums, steadier bass, brighter pad, clearer pluck, thicker vocal. The track count can stay the same while every role becomes more definite.

\begin{figure}[htbp]
  \centering
  \begin{tikzpicture}[font=\small]
    \draw[->] (0,0) -- (10.2,0) node[right] {bars};
    \draw[->] (0,0) -- (0,3.8) node[above] {energy};
    \foreach \x/\lab in {0/1,2/3,4/5,6/7,8/9,10/11,12/13,14/15,16/16} {\draw[gray!25] (\x,0.1) -- (\x,3.5); \node[font=\scriptsize] at (\x,0.05) {\lab};}
    \draw[very thick, yebanblue] (0.0,1.4) -- (2.0,1.6) -- (4.0,1.7) -- (6.0,1.8) -- (8.0,2.7) -- (10.0,2.9) -- (12.0,3.0) -- (14.0,3.1) -- (16.0,3.2);
    \draw[very thick, orange!80] (0.0,1.2) -- (2.0,1.3) -- (4.0,1.5) -- (6.0,1.6) -- (8.0,2.9) -- (10.0,3.1) -- (12.0,3.2) -- (14.0,3.4) -- (16.0,3.5);
    \node[font=\scriptsize, yebanblue, anchor=west] at (16.2,3.2) {folk};
    \node[font=\scriptsize, orange!80, anchor=west] at (16.2,3.5) {pop};
    \node[font=\scriptsize] at (4.0,0.25) {verse};
    \node[font=\scriptsize] at (12.0,0.25) {chorus};
  \end{tikzpicture}
  \caption{Energy curves of the two short songs: folk rises gently, pop lifts more clearly from verse to chorus; neither is a flat line.}
  \label{fig:walkthrough-energy}
\end{figure}

\section{Eight-step walkthrough}
\label{sec:walkthrough-steps}

\subsection*{Step 1: section chart and energy curve}
Write two tables first: a section chart (bars, section, action) and a track chart (track, type, role). Without them, every added track becomes ``add first, decide later''.

\subsection*{Step 2: create and name tracks}
Use the role cards from chapter 1: rhythm, harmony, melody, colour. Name tracks so a human can read them (``Lead'', ``Bass'', ``Pad''), not ``Track 1'' and ``Track 2''.

\subsection*{Step 3: rhythm section}
Build drums and bass first: drums define time and groove, bass defines the root. Solo both and confirm kick and bass root do not fight on the same beat. Folk drums stay light; pop drums can be fuller, but neither should fill every sixteenth.

\subsection*{Step 4: harmony layer}
Add acoustic guitar/keyboard/pad: write the chord movement first, then choose the texture. Folk uses fingerpicking/strumming; pop uses a pad bed. Do not start with five harmony tracks.

\subsection*{Step 5: lead melody and vocal}
Write the lead or record the vocal: make the lyric clear and leave phrase ends open. Folk can rely on the guitar; pop needs double/harmony behind the lead. Never let backing cover the lead.

\subsection*{Step 6: colour layers}
Add strings, woodwinds/brass or electronic pluck by section: enter only when colour/energy is needed, and be able to exit. Both short songs use one or two colours, not every instrument in the book.

\subsection*{Step 7: section-by-section walkthrough}
Solo/Mute through the section chart: does every section enter and exit? Is only one layer leading the low end? Is the lyric clear? Is there a fill/space at the transition? Answer at least once per section; do not just play the whole song once.

\subsection*{Step 8: save and deliver}
Save the project, export WAV and stems, write the receipt. Problems found in step 7 go back to the relevant track; do not force them down with faders.

\section{Walkthrough checklist}
\label{sec:walkthrough-checklist}

\begin{longtable}{p{0.36\textwidth}p{0.56\textwidth}}
\toprule
Check & Pass criterion \\
\midrule
\endfirsthead
\toprule
Check & Pass criterion \\
\midrule
\endhead
Sections & Every section has a clear entrance/exit; none is identical from start to finish \\
Energy & The chorus is more complete, not merely louder \\
Low end & Only one of sub/bass/kick leads; the other two yield \\
Lead vocal & Every word clear; breath at phrase ends; not covered by double/harmony \\
Harmony & Chord movement clear; inner voices do not share the lead's rhythm \\
Colour & Strings/woodwinds/brass/electronic enter only when needed and can exit \\
Transitions & Fill, space, harmonic transition or hard cut between sections, not random switching \\
Project & Track names readable; clips not lost; save succeeds \\
Export & WAV has sample rate/bit depth/duration; stem receipt reports exported/skipped counts \\
\bottomrule
\end{longtable}

\section{Save, export and deliver}
\label{sec:walkthrough-delivery}

\begin{enumerate}
\item \textbf{Save the project}: use File $\rightarrow$ Save project (Ctrl+S) to write back to the same file, or File $\rightarrow$ Save as (Ctrl+Shift+S) to a new name. Saving is an \term{atomic write}{atomic-write}; a power cut halfway does not leave half a file.
\item \textbf{Version history}: use Edit $\rightarrow$ Arrangement version time machine to view the commit chain; click any version node to return to it. This is project editing history, not a mix version.
\item \textbf{Export WAV}: use File $\rightarrow$ Export audio (WAV)... to set output path, sample rate, bit depth, duration and the bus limiter. The receipt reports frames, true peak and limiter work.
\item \textbf{Export stems}: use File $\rightarrow$ Export stems... to export each track separately; the receipt reports how many were exported, skipped and the output directory.
\item \textbf{Do not confuse the diagnostic bundle}: Diagnostics $\rightarrow$ Export project diagnostic report bundle is for troubleshooting, not audio; do not deliver it as a WAV.
\item \textbf{Write the delivery receipt}: record project file, WAV, stem directory, sample rate/bit depth, export time and skipped items. Without a receipt, nobody can reproduce the delivery.
\end{enumerate}

\screenshot{arrangement-walkthrough-folk}{The sixteen-bar folk project: lead, acoustic guitar, bass, light drums, string pad and double, checked section by section}

\screenshot{arrangement-walkthrough-pop}{The sixteen-bar pop project: lead, double, harmony, drums, bass, pad, pluck and brass; chorus layers and the bar-9 brass stab are visible}

\section{Common mistakes and fixes}
\label{sec:walkthrough-mistakes}

\begin{longtable}{p{0.34\textwidth}p{0.58\textwidth}}
\toprule
Common mistake & Fix \\
\midrule
\endfirsthead
\toprule
Common mistake & Fix \\
\midrule
\endhead
Add tracks before planning sections & Write the section chart and track chart first; every track must answer ``what does it do here?'' \\
The chorus is only louder & Make each layer more complete: drums, bass, pad, vocal each add something; do not just raise faders. \\
Both demos use the same track list & Folk is few but clear; pop is each layer clearer. Change the track chart by style, do not copy. \\
Colour layers run throughout & Colour enters at section changes/chorus and can exit. \\
The lead is covered by double/harmony & Lead centred and highest; double same melody, harmony a third below; stack only the final pass. \\
Clips are not checked before saving & Solo/Mute every section before saving; fix muted/lost clips on the spot. \\
Treating the diagnostic bundle as delivery & WAV/stems are the work; the diagnostic bundle is only for troubleshooting. \\
\bottomrule
\end{longtable}

\section{Exercise}
\label{sec:walkthrough-exercise}

\begin{exercise}
Choose folk or pop and build sixteen bars in Yeban:
\begin{enumerate}
\item Write a section chart (eight-bar verse + eight-bar chorus) and a track chart (no more than eight tracks).
\item Follow the eight steps: rhythm $\rightarrow$ harmony $\rightarrow$ lead/vocal $\rightarrow$ colour; solo after adding each track.
\item Tick every item on the walkthrough checklist; fix at least one ``seat-stealing'' problem (low end, lead register or a colour layer that never exits).
\item Save the project, export WAV and stems, and write a delivery receipt; keep the project file, WAV, stem directory and receipt in one folder.
\end{enumerate}
When finished, answer one question: if you could keep only five tracks, which three would you delete? If you can answer clearly, you understand what each layer does.
\end{exercise}

\section{Chapter takeaways}
\label{sec:walkthrough-takeaways}

\begin{itemize}
\item A walkthrough is not ``listen again''; it checks entrances, exits and seat-stealing layer by layer against the section chart.
\item Both short songs share one workflow: section chart $\rightarrow$ track chart $\rightarrow$ rhythm $\rightarrow$ harmony $\rightarrow$ vocal $\rightarrow$ colour $\rightarrow$ check $\rightarrow$ deliver.
\item Folk is few but clear; pop is each layer clearer. Different styles need different track charts.
\item The chorus is ``big'' because every layer is more complete, not because faders or track counts are higher.
\item Saving is atomic; WAV and stems are two deliveries; the diagnostic bundle is not the work.
\item Delivery must include a receipt: project, WAV, stems, parameters and skipped items can all be checked.
\end{itemize}

\begin{helpmapping}
\begin{longtable}{p{0.45\textwidth}p{0.45\textwidth}}
\toprule
Section & Help anchor \\
\midrule
The goal & \texttt{arrangement-walkthrough-goal} \\
Demo one: folk & \texttt{arrangement-folk-demo} \\
Demo two: pop & \texttt{arrangement-pop-demo} \\
Eight-step walkthrough & \texttt{arrangement-walkthrough-steps} \\
Walkthrough checklist & \texttt{arrangement-walkthrough-checklist} \\
Save, export and deliver & \texttt{arrangement-walkthrough-delivery} \\
Common mistakes and fixes & \texttt{arrangement-walkthrough-mistakes} \\
Exercise & \texttt{arrangement-walkthrough-exercise} \\
\bottomrule
\end{longtable}
\end{helpmapping}
